\bookchapter{Κεφάλαιο 48: Γνωρι$\mu$ία $\mu$ε το Android — από το λειτουργικό σύστη$\mu$α στο οικοσύστη$\mu$α ανάπτυξης}

\chapterintro{Συγχαρητήρια για την απόφαση — από «κάποιος που παίζει $\mu$ε το τηλέφωνο» σε «κάποιος που φτιάχνει εφαρ$\mu$ογές τηλεφώνου». Πριν πιάσου$\mu$ε δουλειά, ας απλώσου$\mu$ε τον χάρτη: από πού προήλθε το Android; Γιατί οι developers σε όλο τον κόσ$\mu$ο γράφουν εφαρ$\mu$ογές για αυτό; Τι συ$\mu$βαίνει από τη στιγ$\mu$ή που χτυπάς τα πλήκτρα $\mu$έχρι η εφαρ$\mu$ογή να φτάσει στο τηλέφωνο του χρήστη;}

\booksection{Οι ρίζες του Android: από το «$\mu$ικρό ρο$\mu$ποτάκι» στο «$\mu$εγαλύτερο σύστη$\mu$α του κόσ$\mu$ου»}

Η ιστορία του Android αρχίζει το 2003 — ο $\mu$ηχανικός Andy Rubin ίδρυσε την εταιρεία Android $\mu$ε έναν ταπεινό στόχο: ένα έξυπνο λειτουργικό σύστη$\mu$α που να «καταλαβαίνει την κά$\mu$ερα» καλύτερα. Το 2005 η Google το αγόρασε, το 2007 $\mu$αζί $\mu$ε δεκάδες κατασκευαστές σχη$\mu$άτισαν την Open Handset Alliance (OHA), και τον Σεπτέ$\mu$βριο 2008 κυκλοφόρησε επίση$\mu$α το πρώτο τηλέφωνο Android, το HTC Dream (G1).

Την υπόλοιπη ιστορία τη ξέρου$\mu$ε όλοι: $\mu$έχρι το 2026 το Android είναι το λειτουργικό σύστη$\mu$α κινητών $\mu$ε τη $\mu$εγαλύτερη $\mu$ερίδα αγοράς παγκοσ$\mu$ίως, $\mu$ε δισεκατο$\mu$$\mu$ύρια ενεργές συσκευές. Το τηλέφωνό σου, το TV box σου, η $\mu$εγάλη οθόνη του αυτοκινήτου, το έξυπνο ρολόι σου — πολύ πιθανόν όλα τρέχουν Android. Τι ση$\mu$αίνει αυτό; Ότι κάθε γρα$\mu$$\mu$ή κώδικα που γράφεις έχει την ευ και ρία να τη χρησι$\mu$οποιήσουν οι περισσότεροι άνθρωποι του πλανήτη.

\begin{tcbinfo}{Συ$\mu$βουλή: το DNA ανοιχτού κώδικα του Android}
Ο πυρήνας του Android (AOSP, Android Open Source Project) είναι ανοιχτού κώδικα — κάθε κατασκευαστής $\mu$πορεί να τον πάρει, να τον τροποποιήσει και να τον προσαρ$\mu$όσει. Γι' αυτό η Xiaomi έχει το MIUI (τώρα HyperOS), η Huawei το συ$\mu$βατό στρώ$\mu$α HarmonyOS, η Samsung το One UI. Για τον developer αυτό είναι καλό: $\mu$αθαίνεις «ανάπτυξη Android», όχι «ανάπτυξη για ένα brand» — $\mu$ια δουλειά που δουλεύει παντού.

\end{tcbinfo}

Και η εξέλιξη των εκδόσεων έχει ενδιαφέρον: από το 1.5 $\mu$ε το Cupcake, κάθε $\mu$εγάλη έκδοση έπαιρνε όνο$\mu$α γλυκού — Donut, Eclair, Frozen Yogurt, Gingerbread… και $\mu$όνο $\mu$ε το Android 10, το 2019, πέρασε σε σκέτα νού$\mu$ερα. Μέχρι τον Αύγουστο 2026, η νεότερη επίση$\mu$η έκδοση είναι το Android 17· σε επίπεδο ανάπτυξης, το «ελάχιστο υποστηριζό$\mu$ενο» που πραγ$\mu$ατικά σε αφορά, στις mainstream εφαρ$\mu$ογές ορίζεται συνήθως ανά$\mu$εσα στο Android 8 (API 26) και στο Android 10 (API 29) — όσο υποστηρίζεις συσκευές λίγων ετών, τόσο $\mu$εγαλύτερο είναι το κοινό σου.

\begin{center}
\begin{tabularx}{\linewidth}{|X|X|X|}
\hline
\textbf{Έκδοση} & \textbf{Κωδικό όνο$\mu$α / Έτος} & \textbf{Τα ση$\mu$αντικά για τον developer} \tabularnewline
\hline
Android 8 (API 26) & Oreo · 2017 & Κανάλια ειδοποιήσεων, περιορισ$\mu$οί στο background — το ελάχιστο όριο για πολλές παλιές συσκευές \tabularnewline
Android 10 (API 29) & 2019 & Έρχεται το scoped storage· οι κανόνες πρόσβασης σε αρχεία αλλάζουν ριζικά \tabularnewline
Android 12 (API 31) & 2021 & Material You $\mu$ε δυνα$\mu$ικά χρώ$\mu$ατα· $\mu$εγάλη αναβάθ$\mu$ιση στην όψη του UI \tabularnewline
Android 13 (API 33) & 2022 & Άδεια ειδοποιήσεων κατά την εκτέλεση· παράθυρο που ρωτά «επιτρέπεις τις ειδοποιήσεις;» \tabularnewline
Android 14-15 (API 34-35) & 2023-2024 & Προβλεπτικό back, σελίδες 16KB, περαιτέρω σύσφιξη ιδιωτικότητας \tabularnewline
Android 16-17 (API 36-37) & 2025-2026 & Πιο γρήγορος ρυθ$\mu$ός $\mu$ικρών εκδόσεων, ενσω$\mu$άτωση AI (π.χ. επεκτάσεις Gemini) \tabularnewline
\hline
\end{tabularx}
\end{center}

\begin{center}\small\color{gray}Σύντο$\mu$η ιστορία εκδόσεων του Android (κατά τον Αύγουστο 2026, τα ση$\mu$εία-κλειδιά για ανάπτυξη)\end{center}

\booksection{Σε ποια γλώσσα γράφεις: Java ή Kotlin;}

Στην ιστορία του Android ε$\mu$φανίστηκαν δύο κύριες γλώσσες. Η Java είναι ο «παλιός στρατιώτης της πρώτης ώρας» — την είχε το Android από τη γέννησή του· το Kotlin ανακηρύχθηκε το 2017 από τη Google «γλώσσα πρώτης κατηγορίας» και από το 2019 η επίση$\mu$η γρα$\mu$$\mu$ή είναι σαφές «Kotlin-first». Το 2026, σχεδόν όλα τα νέα projects είναι σε Kotlin και τα παλιά projects $\mu$εταφέρονται σταδιακά.

Πά$\mu$ε στα ίσια: αν είσαι εντελώς αρχάριος, $\mu$άθε ευθεία Kotlin και $\mu$η γυρίζεις πίσω στη Java. Τρεις λόγοι: πρώτον, ο κώδικας Kotlin είναι 30\%-40\% $\mu$ικρότερος από τη Java (ίδια λειτουργία, πιο σύντο$\mu$ος και καθαρός)· δεύτερον, προστατεύει εκ φύσεως από το «null pointer», τον εφιάλτη των αρχάριων· τρίτον, όλα τα σύγχρονα $\mu$αθή$\mu$ατα, τα επίση$\mu$α έγγραφα και τα εργαλεία προγρα$\mu$$\mu$ατισ$\mu$ού $\mu$ε AI (θυ$\mu$άσαι το βιβλίο $\mu$ας για τον προγρα$\mu$$\mu$ατισ$\mu$ό $\mu$ε AI;) βασίζονται στο Kotlin. Java θα συναντήσεις $\mu$όνο όταν συντηρείς παλιά projects ή διαβάζεις παλιό κώδικα.

Στυλ Java ($\mu$ακροσκελές)

\begin{lstlisting}
public class MainActivity extends Activity \{
@Override
protected void onCreate(Bundle b) \{
super.onCreate(b);
setContentView(R.layout.activity);
\}
\}
\end{lstlisting}

Στυλ Kotlin (καθαρό)

\begin{lstlisting}
class MainActivity : ComponentActivity() \{
override fun onCreate(b: Bundle?) \{
super.onCreate(b)
setContent \{ App() \}
\}
\}
\end{lstlisting}

Για το ίδιο πράγ$\mu$α, το Kotlin γλιτώνει σωρό άγκιστρα και boilerplate. Στο κεφάλαιο 50 θα το $\mu$άθου$\mu$ε σωστά· για την ώρα κράτα ένα: η «κοινή γλώσσα» της ανάπτυξης Android λέγεται Kotlin.

\booksection{Πώς γεννιέται $\mu$ια εφαρ$\mu$ογή;}

Πολλοί νο$\mu$ίζουν ότι «η ανάπτυξη εφαρ$\mu$ογής» ση$\mu$αίνει ότι ανοίγεις τον υπολογιστή και γράφεις αδιάκοπα κώδικα. Στην πραγ$\mu$ατικότητα, από την ιδέα $\mu$έχρι το τηλέφωνο του χρήστη, $\mu$ια εφαρ$\mu$ογή περνά περίπου από αυτή τη γρα$\mu$$\mu$ή παραγωγής:

Ιδέα

Σχεδίαση (οθόνη+λειτουργίες)

Γράφει κώδικα

Μεταγλώττιση και πακετάρισ$\mu$α (APK/AAB)

Τεστ (emulator+πραγ$\mu$ατική συσκευή)

Δη$\mu$οσίευση (κατάστη$\mu$α/πακέτο εγκατάστασης)

\begin{center}\small\color{gray}Σχή$\mu$α 1: Ολόκληρη η ροή $\mu$ιας Android εφαρ$\mu$ογής, από την ιδέα $\mu$έχρι το κατάστη$\mu$α\end{center}

Στη ροή υπάρχουν δύο έννοιες που πρέπει να γνωρίζεις από τώρα:

\begin{itemize}
  \item APK: η $\mu$ορφή του πακέτου εγκατάστασης Android εφαρ$\mu$ογών (.apk)· $\mu$ε διπλό κλικ εγκαθίσταται στο τηλέφωνο, όπως το .exe στα Windows. Για τις προσωπικές σου δοκι$\mu$ές και για να το $\mu$οιράζεσαι σε φίλους για πείραγ$\mu$α, αυτό χρησι$\mu$οποιείς.
  \item AAB: Android App Bundle (.aab), η «προηγ$\mu$ένη $\mu$ορφή πακεταρίσ$\mu$ατος» που απαιτείται για το Google Play. Το κατάστη$\mu$α το «κόβει» δυνα$\mu$ικά ανάλογα $\mu$ε το $\mu$οντέλο του τηλεφώνου του χρήστη, ώστε το πακέτο που κατεβαίνει να είναι $\mu$ικρότερο. Μόνος σου δεν το εγκαθιστάς — πηγαίνει $\mu$όνο στο κατάστη$\mu$α.
\end{itemize}

Ακό$\mu$α ένας όρος που θα ακούς συνέχεια: SDK (Software Development Kit, κιτ ανάπτυξης λογισ$\mu$ικού). Είναι ολόκληρος ο «εξοπλισ$\mu$ός» που προσφέρει επίση$\mu$α το Android — τις βιβλιοθήκες API κάθε έκδοσης, τα εργαλεία debugging, τις εικόνες του emulator κ.ά. Κλάσεις όπως το Activity και το Button που καλείς στον κώδικα προέρχονται από το SDK.

\booksection{Το τεχνολογικό τοπίο του Android σε $\mu$ια $\mu$ατιά: η «λίστα οπλισ$\mu$ού» σου}

Πριν ξεκινήσεις σοβαρά, γνώρισε την τυπική δια$\mu$όρφωση των mainstream Android projects του 2026 — αν δεν την καταλαβαίνεις τώρα, δεν πειράζει· στο τέλος του βιβλίου θα τη ξαναδείς και όλα θα σου φαίνονται καθαρά:

\begin{center}
\begin{tabularx}{\linewidth}{|X|X|X|}
\hline
\textbf{Επίπεδο} & \textbf{Σύγχρονη επιλογή (2026)} & \textbf{Σε τι χρησι$\mu$εύει} \tabularnewline
\hline
Γλώσσα & Kotlin 2.x & Η γλώσσα για τη λογική \tabularnewline
Διεπαφή & Jetpack Compose + Material 3 & Σχεδιάζεις τη διεπαφή $\mu$ε κώδικα (η κύρια γρα$\mu$$\mu$ή του βιβλίου) \tabularnewline
Αρχιτεκτονική & MVVM + ViewModel & Ξεχωρίζεις διεπαφή και λογική, ώστε ο κώδικας να $\mu$ην $\mu$περδεύεται \tabularnewline
Αποθήκευση δεδο$\mu$ένων & Room (βάση δεδο$\mu$ένων) + DataStore (ρυθ$\mu$ίσεις) & Αποθηκεύεις δεδο$\mu$ένα στο τηλέφωνο \tabularnewline
Δίκτυο & Retrofit + OkHttp + Kotlin Coroutines & Μιλάς $\mu$ε τον server \tabularnewline
Εργαλείο build & Gradle + Kotlin DSL & Μετατρέπεις τον κώδικα σε πακέτο εγκατάστασης \tabularnewline
Περιβάλλον ανάπτυξης & Android Studio & Το «όλα-σε-ένα» IDE της Google \tabularnewline
\hline
\end{tabularx}
\end{center}

\begin{center}\small\color{gray}Ο mainstream tech stack του Android το 2026, σε $\mu$ια $\mu$ατιά\end{center}

Μη τρο$\mu$άζεις $\mu$ε τον σωρό των αγγλικών όρων. Είναι όπως τα σκεύη της κουζίνας — για την ώρα αρκεί να ξέρεις ότι «για να $\mu$αγειρέψεις, $\mu$παίνεις στην κουζίνα»· πώς χρησι$\mu$οποιείται το κάθε εργαλείο, θα σου το ξετυλίγου$\mu$ε ένα-ένα στα επό$\mu$ενα κεφάλαια. Και το καλό νέο: όταν το Android Studio δη$\mu$ιουργεί νέο project, σετάρει όλο αυτό το σετ αυτό$\mu$ατα — εσύ απλώς $\mu$αγειρεύεις.

\booksection{Android vs iOS: το πρώτο $\mu$άθη$\mu$α «επιλογής πλευράς» για τον developer}

Ίσως να έχεις στα χέρια σου ταυτόχρονα συσκευή Android και συσκευή Apple. Ως developer, αξίζει να γνωρίζεις από τώρα τις διαφορές των δύο πλευρών — δεν πρόκειται για το «ποιο είναι καλύτερο», αλλά για τον ρεαλιστικό υπολογισ$\mu$ό «από ποια πλευρά σου αξίζει να ξεκινήσεις»:

\begin{center}
\begin{tabularx}{\linewidth}{|X|X|X|}
\hline
\textbf{Κριτήριο} & \textbf{Ανάπτυξη Android} & \textbf{Ανάπτυξη iOS} \tabularnewline
\hline
Υπολογιστής ανάπτυξης & Windows / Mac / Linux — όλα ισχύουν & Μόνο Mac (αποκλειστικό Xcode) \tabularnewline
Εργαλεία ανάπτυξης & Android Studio (δωρεάν) & Xcode (δωρεάν) + ετήσιο 99\$ για δη$\mu$οσίευση \tabularnewline
Έλεγχος δη$\mu$οσίευσης & Χαλαρός, από λίγες ώρες έως λίγες $\mu$έρες & Αυστηρός, 1-7 η$\mu$έρες είναι ο κανόνας \tabularnewline
Εγκατάσταση και δοκι$\mu$ή & APK: το $\mu$οιράζεσαι και το εγκαθιστάς ελεύθερα & Περιορισ$\mu$ένη διανο$\mu$ή· για τεστ χρειάζεται υπογραφή/TestFlight \tabularnewline
Περιβάλλον εσόδων & Κυρίως διαφη$\mu$ίσεις / τοπικά κανάλια & Καλύτερες συνήθειες πληρω$\mu$ών εντός εφαρ$\mu$ογής στο App Store \tabularnewline
Κόστος εκ$\mu$άθησης & Πολλές έννοιες (τα τέσσερα βασικά components κ.ά.) & Κλειστό οικοσύστη$\mu$α αλλά ενιαίο, $\mu$ε άριστη τεκ$\mu$ηρίωση \tabularnewline
\hline
\end{tabularx}
\end{center}

\begin{center}\small\color{gray}Σύγκριση των διαδρο$\mu$ών ανάπτυξης Android και iOS (Αύγουστος 2026)\end{center}

Πά$\mu$ε στα ίσια: αν ο υπολογιστής σου δεν είναι Mac, το ερώτη$\mu$α δεν χρειάζεται καν σκέψη — το Android είναι η $\mu$οναδική επιλογή. Αν έχεις Mac, τα $\mu$έρη 8 και 9 του βιβλίου $\mu$πορούν να διαβαστούν παράλληλα· η λογική του προγρα$\mu$$\mu$ατισ$\mu$ού είναι κοινή, και ό, τι $\mu$άθεις στη $\mu$ία πλευρά θα σε βοηθήσει να $\mu$άθεις την άλλη πολύ πιο γρήγορα.

\begin{tcbpractice}{Σειρά σου: δες το «Android» $\mu$έσα στο τηλέφωνό σου}
Άνοιξε στο τηλέφωνο [Ρυθ$\mu$ίσεις] [Σχετικά $\mu$ε το τηλέφωνο] (ή «Σχετικά $\mu$ε τη συσκευή»), βρες τη γρα$\mu$$\mu$ή «Έκδοση Android» — αν την πατήσεις αρκετές φορές συνεχό$\mu$ενα, θα ενεργοποιηθεί και το $\mu$ίνι παιχνίδι-έκπληξη του συστή$\mu$ατος. Ση$\mu$είωσε τον αριθ$\mu$ό της έκδοσής σου και σύγκρινε τον $\mu$ε τον παραπάνω πίνακα για να δεις σε ποια γενιά ανήκει. Το τηλέφωνό σου είναι η πρώτη σου «συσκευή δοκι$\mu$ών».

\end{tcbpractice}

\begin{tcbsummary}{Σύνοψη κεφαλαίου}
\begin{itemize}
  \item Το Android γεννήθηκε το 2008· το 2026 είναι το λειτουργικό κινητών $\mu$ε τη $\mu$εγαλύτερη $\mu$ερίδα αγοράς, $\mu$ε δισεκατο$\mu$$\mu$ύρια ενεργές συσκευές
  \item Νεότερη επίση$\mu$η έκδοση: Android 17· οι mainstream εφαρ$\mu$ογές στηρίζουν τουλάχιστον Android 8-10, για να είναι $\mu$εγάλο το κοινό
  \item Για νέα projects: πάντα Kotlin — 30\%+ λιγότερος κώδικας, προστασία από null, επίση$\mu$η προτί$\mu$ηση της Google
  \item APK για προσωπική εγκατάσταση, AAB για το κατάστη$\mu$α· το SDK είναι το επίση$\mu$ο «κιτ εξοπλισ$\mu$ού» ανάπτυξης
  \item Τυπικό tech stack του 2026: Kotlin + Compose + MVVM + Room + Retrofit + Gradle + Android Studio
\end{itemize}

\end{tcbsummary}
